%! TEX root = master.tex
\chapter{1}{Statements and Proof}

The first kind of proof is the lowly conditional, $P \implies Q$, read
``$P$ implies $Q$'' or ``if $P$, then $Q$''. To prove it, we take what we are
given in $P$, together with what we already know to be true, and use it to
reach $Q$. This chapter fixes what a conditional means and what may be
traded for it.

\section{The conditional}\label{sec:conditional}

\begin{definition}[Conditional]\label{def:conditional}
  For statements $P$ and $Q$, the conditional $P \implies Q$ is false when
  $P$ is true and $Q$ is false, and true in every other case. $P$ is the
  \emph{hypothesis} and $Q$ the \emph{conclusion}. Three related statements
  are
  \begin{center}
    \begin{tabular}{ll}
      converse & $Q \implies P$, \\
      inverse & $\lnot P \implies \lnot Q$, \\
      contrapositive & $\lnot Q \implies \lnot P$.
    \end{tabular}
  \end{center}
\end{definition}

\begin{remark}
  A conditional with a false hypothesis is true, however odd its
  conclusion. ``If $1 = 2$, then the moon is cheese'' is a true statement.
  A conditional promises nothing when its hypothesis fails.
\end{remark}

%! TEX root = ../master.tex
\begin{theorem}[Contrapositive]\label{thm:contrapositive}
  For any statements $P$ and $Q$, the conditional $P \implies Q$ and its
  contrapositive $\lnot Q \implies \lnot P$ are logically equivalent: they are
  true in exactly the same cases.
\end{theorem}

\begin{proof}
  Compare the two in every case:
  \[
    \begin{array}{cccc|c|c}
      P & Q & \lnot P & \lnot Q & P \implies Q & \lnot Q \implies \lnot P \\
      \hline
      T & T & F & F & T & T \\
      T & F & F & T & F & F \\
      F & T & T & F & T & T \\
      F & F & T & T & T & T
    \end{array}
  \]
  The last two columns agree in every row.
\end{proof}


\begin{remark}
  The converse is \emph{not} equivalent to the conditional. In the second
  row of the table ($P$ true, $Q$ false), $P \implies Q$ is false but
  $Q \implies P$ is true. The inverse is the contrapositive of the converse,
  so it stands or falls with the converse.
\end{remark}

\section{Three ways to prove \texorpdfstring{$P \implies Q$}{P implies Q}}\label{sec:methods}

\begin{description}[font=\normalfont\scshape]
  \item[Direct.] Assume $P$ and deduce $Q$ step by step.
  \item[Contrapositive.] Assume $\lnot Q$ and deduce $\lnot P$. By
    \cref{thm:contrapositive}, this proves $P \implies Q$.
  \item[Contradiction.] Assume $P$ and $\lnot Q$ together, and deduce
    something false. Then the case ``$P$ true, $Q$ false'' cannot happen,
    and that is the only case in which $P \implies Q$ fails.
\end{description}

Each method appears in the next two chapters. The proof that $n^2$ even
implies $n$ even (\cref{thm:square-parity}) goes by contrapositive, and the
irrationality of $\sqrt2$ (\cref{thm:sqrt2-irrational}) by contradiction.

\begin{exercise}\label{ex:converse-square}
  Let $n$ be an integer. State the converse, inverse and contrapositive of
  ``if $n$ is even, then $n^2$ is even'', and decide which of the four
  statements are true.
\end{exercise}

\begin{solution}
  \begin{itemize}
    \item Converse: if $n^2$ is even, then $n$ is even.
    \item Inverse: if $n$ is odd, then $n^2$ is odd.
    \item Contrapositive: if $n^2$ is odd, then $n$ is odd.
  \end{itemize}
  Here ``not even'' is the same as ``odd'' by \cref{thm:parity}. All four
  statements are true: the original and its contrapositive by
  \cref{thm:parity-arithmetic}, and the converse and inverse by
  \cref{thm:square-parity}. This is unusual. A conditional and its converse
  both hold exactly when $P$ and $Q$ are equivalent, as they are here.
\end{solution}
